% This section follows the shared normalization in (6).
\section{Consecutive cubic rows and the first quartic}
\label{sec:cubic}

Let \(n=3j+m\), where \(0\le m<j\). In the coordinate \(y=D(t-1)\) of
Section~5, let \(P\) and \(Q\) denote the two rows divided by their own
values at \(t=1\). Write
\[
 P=1+ay+by^2+cy^3,\qquad Q-P=y(d_0+e_0y+f_0y^2+zy^3).
\]
By (6), \(c>0\), with \(c=K_0^4/K_{3j}\) at \(m=0\).
The coefficient \(z\) vanishes for \(m<j-1\); at \(m=j-1\),
\(z=K_0^5/K_{4j}>0\). Put \(\kappa=z/c\) and
\[
 (d,e,f)=(d_0,e_0,f_0)-\kappa(1,a,b).
\]
Then \(Q=(1+\kappa y)P+yH\) with \(H=d+ey+fy^2\).

\subsection{A sign criterion using the two separator roots}

Suppose \(a,b,c,d,e,f>0\) and \(h=df/e^2\le4/25\).
The number \(v=2/(1+\sqrt{1-4h})\) lies in \([1,5/4]\) and satisfies
\(1/v=1-hv\). The two roots of \(H\) are \(-W,-w\), where
\(w=dv/e\), \(W=e/(fv)\), and \(w/W=hv^2=v-1\le1/4\).
For any upper bound \(V\ge v\) we claim that the tests
\begin{equation}
 ad>e,\qquad
 (ad-e)e>Vd\max(0,bd-f),\qquad
 V\frac{af}{be}+\frac{ce}{bf}<1
 \tag{7}
 \label{eq:cubic-tests}
\end{equation}
give \(P(-w)<0<P(-W)\). For the first sign, drop the negative cubic term
and use \(1/v=1-hv\) to obtain
\[
 \frac{P(-w)}{v}\le 1-\frac{ad}{e}
                   +\frac{dv(bd-f)}{e^2}<0.
\]
The far-point evaluation is
\[
 P(-W)=1+bW^2
       \left(1-v\frac{af}{be}-\frac{1}{v}\frac{ce}{bf}\right)>0.
\]
Here \(v\ge1\) and \(v\le V\) allow the last comparison in
\eqref{eq:cubic-tests}.

Since \(c>0\) and \(P(0)=1\), the signs at \(-\infty,-W,-w,0\)
give precisely three simple roots \(r_1<-W<r_2<-w<r_3<0\).
At these roots the evaluations of \(Q=yH\) have signs \(-,+,-\),
respectively. Together with \(Q(0)=1\), this isolates three new roots
to the right of \(r_1\), one in each interval \((r_1,r_2)\),
\((r_2,r_3)\), and \((r_3,0)\). When \(z>0\), the positive quartic
leading term gives one more root to the left of \(r_1\).
The degree exhausts the roots in both cases. This is the strict
root transfer of Liu--Wang \cite[Theorem~2.1]{lw}, with separator \(H\)
and multiplier \(y<0\) at the old roots. The increasing change back
to \(t\) preserves order. The original coefficients are nonnegative,
with positive constant, so every real root has \(t<0\).

\subsection{Coefficient estimates}

Keep \(K\) and \(B\) from Section~5. Write \(K_i=k_i=[x^i]K(x)\).
The general formula (6) gives every \(y^h\) coefficient. Define
\(g_i=[x^i]K^2\), \(w_i=[x^i]K^3\), \(v_i=[x^i]K^4\), with zero values
at negative indices, and write \(\Delta g_i=g_{i+1}-g_i\) and likewise
for the other kernels. We use the explicit expressions for \(g_i,w_i\)
from Sections~2 and~5. In particular, at
\(i\ge0\), \(\Delta w_i\) differs from \(i+2+3\rho/2\) by at most
\(0.48\,2^{-i}\).

Twice differentiating \(K'=K^2-\rho K\) yields
\[
 6K^4=K'''+6\rho K''+11\rho^2K'+6\rho^3K .
\]
Consequently for \(i\ge0\),
\[
 v_i=\tfrac16(i+1)(i+2)(i+3)K_{i+3}
      +\rho(i+1)(i+2)K_{i+2}
      +\tfrac{11}{6}\rho^2(i+1)K_{i+1}+\rho^3K_i.
\]
By (1), replacing the four \(K\)'s by 1 costs at most
\(1/1024+0.7/64+(11/6)(0.7)^2/16+0.7^3/2<0.24\).
Here \((i+1)(i+2)(i+3)\le6\cdot8^i\),
\((i+1)(i+2)\le2\cdot4^i\), and \(i+1\le2^i\).
The forward difference of the polynomial part is
\((i+2)(i+3)/2+2\rho(i+2)+11\rho^2/6\). Using \(0.69<\rho<0.70\)
and a difference error at most \(0.48\), we obtain
\begin{align*}
 i^3/6 &\le v_i\le L_3(i):=i^3/6+1.70i^2+4.84i+3.89,\\
 \ell_3(i):=i^2/2+3.87i+6 &\le \Delta v_i
                         \le U_3(i):=i^2/2+3.90i+7.30 .
\end{align*}
The additional nonzero increments at index \(-1\) are \(K_0^2<2\),
\(K_0^3<3\), and \(K_0^4<4\), respectively.

For \(j\ge64\), let \(q=1/j\), \(s=j+m\), \(t_*=2j+m\), and set
\[
 M=(1-0.7q)^{-1},\qquad \eta=q/100,\qquad
 \beta=0.69q/(1+q).
\]
The coefficients \(B_i\) are nonnegative, have sum at most \(M\), and
satisfy \(B_0=1,\ B_1\ge\beta\), because \(B_i\le (0.7q)^i\).
At rank \(N=n\) or \(n+1\), each numerator of order \(1,2,3\) in (6)
is at most \(2(N+7)^3\) by the preceding estimates
and \(M^3<2\). Dividing by \(K_N\) changes it by at most \(\eta\):
\[
 |K_N^{-1}-1|\le8^{-N},\qquad
 200N(N+7)^3\le8^N,\qquad \frac{1}{100N}\le\eta .
\]
The middle bound holds at \(N=100\), and the ratio of the polynomial
expression from \(N\) to \(N+1\) is at most \((101/100)^4<8\)
for all \(N\ge100\). An increment therefore changes by at most \(2\eta\).

The resulting unscaled bounds are as follows:
\begin{equation}
\begin{aligned}
 t_*+1.68+\beta t_*-\eta &\le a \le M(t_*+1.72)+\eta,\\
 1+\beta-3\eta &\le d_0\le1+2(M-1)-\beta+3\eta,\\
 s^2/2+2.52s+2+\beta s^2-\eta &\le b
                   \le M^2(s^2/2+2.55s+2.86)+\eta,\\
 s+3+2\beta s-2\eta &\le e_0\le M^2(s+3.54)+2\eta,\\
 m^3/6-\eta &\le c \le M^3 L_3(m)+\eta,\\
 \ell_3(m)-2\eta &\le f_0\le M^3 U_3(m)+2\eta .
\end{aligned}
\tag{8}
\label{eq:cubic-bounds}
\end{equation}
For \(a\), retain \(g_{t_*}+B_1g_{t_*-1}\), with \(g_{t_*}\ge t_*+1.68\)
and \(g_{t_*-1}\ge t_*\). The unshifted term is at most \(t_*+1.72\)
and every shifted term is also at most this bound. For \(d_0\), the
unshifted and once-shifted increments differ from 1 by at most
\(0.525/4^{j-1}\le\eta/2\). Their sum, with weights \(1,B_1\), differs
from \(1+B_1\) by at most \(\eta\). The remaining mass is at most
\(M-1-B_1\), and the increments are nonnegative and at most 2.
For the lower bound on \(b\), retain \(w_s+2B_1w_{s-1}\), using
\(w_s\ge s^2/2+2.52s+2\) and \(w_{s-1}\ge s^2/2\).
For \(e_0\), the corresponding increments are bounded below by
\(\Delta w_s\ge s+3\) and \(\Delta w_{s-1}\ge s\), since \(s\ge64\).
The upper bounds use the polynomial or increment upper bound at \(s\),
which also bounds all shifted values, and the total mass \(M^2\).
The last two lines follow from the bounds for \(v_i,\Delta v_i\),
retaining the unshifted term for the lower bounds and using \(M^3\)
for the upper bounds. The extra increments at index \(-1\) respect
the same upper bounds.

In an ordinary pair, \(\kappa=0\) and \(f=f_0\ge\ell_3(m)-2\eta>0\).
To bound the actual ratio \(c/f_0\) using the common \(m\), use
\begin{equation}
 M^3L_3(m)+\eta
 \le \tfrac13 M^3(m+3)\bigl(\ell_3(m)-2\eta\bigr).
 \tag{9}
 \label{eq:cubic-ratio}
\end{equation}
Indeed \(\eta\le1/6400\) and \(M^3\ge1\).
After these substitutions, the difference between the right-hand and
left-hand sides of \eqref{eq:cubic-ratio}, multiplied by \(3/M^3\),
is bounded below by
\[
 \frac{40503}{6400}+\frac{9887}{3200}m+\frac{27}{100}m^2>0.
\]
Hence \(c/f_0\le M^3(m+3)/3\).
At the degree increase \(m=j-1\), we have \(0<z<6\).
The bound \(c\ge(j-1)^3/6-\eta>6j^3/50\) follows
already at \(j=64\) after division by \(j^3\), with the lower bound
increasing in \(j\). Thus \(0<\kappa<50/j^3\).

\subsection{The rational boxes}

Use the units
\[
 \alpha=a/j,\quad \beta_2=b/j^2,\quad \gamma=c/j^3,\quad
 d,\quad E=e/j,\quad F=f/j^2 .
\]
The comparisons in \eqref{eq:cubic-tests} and \(h=dF/E^2\) are invariant
under this change. Partition the rectangle
\((q,m/j)\in[0,1/64]\times[0,1]\) into \(6\cdot128\) closed cells, with
intervals \([q_0,q_1]\) and \([u_0,u_1]\).
Evaluate \eqref{eq:cubic-bounds} after substituting \(q=1/j\) and
\(u=m/j\) and passing to these units. Each positive monomial takes its
lower or upper endpoint as appropriate, and each subtracted quantity
takes the opposite endpoint. The functions \(M\), \(\eta\), and
\(\beta=0.69q/(1+q)\) are increasing in \(q\).
For example, with \(M_1=M(q_1)\) and \(\eta_1=\eta(q_1)\), the bounds
on \(\beta_2\) are
\[
\begin{aligned}
\underline{\beta_2}
 &= (1+u_0)^2/2+2.52(1+u_0)q_0+2q_0^2
           +\beta(q_0)(1+u_0)^2-\eta_1q_1^2,\\
\overline{\beta_2}
 &= M_1^2\bigl((1+u_1)^2/2+2.55(1+u_1)q_1+2.86q_1^2\bigr)
           +\eta_1q_1^2 .
\end{aligned}
\]
A cell can contain the quartic entry exactly when
\(q_0+u_0\le1\le q_1+u_1\). The program checks that case in every such
cell, reducing the lower bounds of \(d,E,F\) by
\(50q_1^3(1,\overline{\alpha},\overline{\beta_2})\), respectively.
The upper bounds are unchanged. For an ordinary cell take
\(G=M_1^3(u_1+3q_1)/3\) as an upper bound on \(\gamma/F\) by
\eqref{eq:cubic-ratio}. At a possible quartic entry, the corrected lower
bound \(\underline F\) is positive and we take \(G=\overline\gamma/\underline F\).
The actual ordinary \(F\) is positive even in a cell with nonpositive
\(\underline F\). The ratio estimate avoids dividing by this lower
bound and covers \(m=0\). The edge \(q=0\) only closes the grid, since
all finite cutoffs have \(q>0\).

For a fixed cell let \(h_*=\overline d\,\overline F/\underline E^2\).
Starting at \(V=5/4\), update \(V\gets(1-h_*V)^{-1}\) five times.
If \(v\le V\) and \(h\le h_*\), then \(0<1-h_*V\le1-hv\), so the
new \(V\) remains an upper bound for \(v\). When \(h_*<4/25\),
all iterates are in \([1,5/4]\) and the denominators are positive.
With \(A_*=\underline\alpha\,\underline d-\overline E\), the three
box margins for \eqref{eq:cubic-tests} are
\[
 A_*,\quad
 A_*\underline E-
 V\overline d\,\max(0,\overline{\beta_2}\,\overline d-\underline F),
 \quad
 1-\frac{V\overline\alpha\,\overline F}{\underline{\beta_2}\,\underline E}
   -\frac{G\overline E}{\underline{\beta_2}}.
\]
The first margin is positive before the product by \(\underline E\)
is used. The exact computation checks these margins, the required
positive lower bounds, and \(h_*<4/25\). The 768 ordinary checks and
nine possible quartic checks have the following lower bounds on
their minima, rounded down:
\begin{center}
\begin{tabular}{c|rrrr}
 & \(4/25-h_*\) & \(A_*\) & Near margin & Far margin\\
\hline
777 checks & \(0.021951\) & \(0.964572\) & \(0.035763\) & \(0.136898\)
\end{tabular}
\end{center}
The attached checker uses rational arithmetic for every comparison.
Each actual point in the stated range belongs to a closed cell with
the required flag. This proves the sign tests for every \(j\ge64\).

\subsection{Finite complement}

For \(2\le j\le63\), use the exact last-block recurrence of Section~4,
and compare binomial moments with the independent coefficient
extraction of Section~3. In the original variable, let \(R\) be a positive
multiple of the negative remainder of \(Q\) on division by \(P\),
then similarly \(S=-\operatorname{rem}(P,R)\) and
\(T=-\operatorname{rem}(R,S)\) up to positive factors.
Integer pseudo-division uses positive leading coefficients, so the
scale factors preserve signs. The attached finite checker gives
degrees \(2,1,0\) and positive leading coefficients for \(R,S,T\)
on every one of the \(2015=\sum_{j=2}^{63}j\) pairs, including
all 62 quartic entries.

At the root of \(S\), \(R\) is negative because \(T\) is positive.
Thus the quadratic \(R\) has two distinct real roots straddling that
point. At its roots \(P\) has signs \(+,-\), since there \(P\) has the
opposite sign of \(S\). The positive leading coefficient of the old
cubic gives one old root in each interval split by the two roots of
\(R\). At the three old roots, \(Q\) has signs \(-,+,-\).
All original coefficients are nonnegative and the constant is positive,
so all old roots are negative and the upper constant gives the
rightmost new root before zero. The quartic leading term gives
the additional leftmost root when present. The degree count again
proves simplicity and completeness. This completes Theorem~2.
